% ============================================================
%  Section 2: Background
% ============================================================

\section{Background}

\newcommand{\codecell}[2]{%
  {\fontsize{17.5}{19}\selectfont\sffamily\bfseries\textcolor{mDarkTeal}{#1:}}%
  \hspace{0.08em}{\fontsize{7.8}{9.4}\selectfont\ttfamily\bfseries\textcolor{black!55}{#2}}%
}

\begin{frame}{Fixed-Length Codes: ASCII}
  \centering
  \resizebox{0.94\textwidth}{!}{%
  \begin{tikzpicture}
    \node[
      fill=mDarkTeal!3,
      draw=mDarkTeal!14,
      rounded corners=2pt,
      inner xsep=9pt,
      inner ysep=9pt
    ] {
      \begin{minipage}{0.98\textwidth}
        \centering
        {\fontsize{25}{27}\selectfont\sffamily\bfseries
          \textcolor{mDarkTeal}{ASCII $\longrightarrow$ }%
          \textcolor{mLightBrown}{8.0}%
          \hspace{0.18em}%
          \textcolor{mDarkTeal}{ bits/ch}}

        \vspace{0.75em}
        \setlength{\tabcolsep}{0pt}
        \renewcommand{\arraystretch}{1.95}
        \begin{tabular*}{\textwidth}{@{\extracolsep{\fill}}llllll}
          \codecell{a}{01100001} & \codecell{b}{01100010} & \codecell{c}{01100011} & \codecell{d}{01100100} & \codecell{e}{01100101} & \codecell{f}{01100110} \\
          \codecell{g}{01100111} & \codecell{h}{01101000} & \codecell{i}{01101001} & \codecell{j}{01101010} & \codecell{k}{01101011} & \codecell{l}{01101100} \\
          \codecell{m}{01101101} & \codecell{n}{01101110} & \codecell{o}{01101111} & \codecell{p}{01110000} & \codecell{q}{01110001} & \codecell{r}{01110010} \\
          \codecell{s}{01110011} & \codecell{t}{01110100} & \codecell{u}{01110101} & \codecell{v}{01110110} & \codecell{w}{01110111} & \codecell{x}{01111000} \\
          \codecell{y}{01111001} & \codecell{z}{01111010} & \codecell{0}{00110000} & \codecell{1}{00110001} & \codecell{2}{00110010} & \codecell{3}{00110011} \\
          \codecell{4}{00110100} & \codecell{5}{00110101} & \codecell{6}{00110110} & \codecell{7}{00110111} & \codecell{8}{00111000} & \codecell{9}{00111001}
        \end{tabular*}
      \end{minipage}
    };
  \end{tikzpicture}%
  }

  \vspace{0.2em}
  \small Every character gets the same 8-bit budget, regardless of how often it appears.
\end{frame}

\begin{frame}{Variable-Length Codes: Huffman}
  \centering
  \resizebox{0.94\textwidth}{!}{%
  \begin{tikzpicture}
    \node[
      fill=mDarkTeal!3,
      draw=mDarkTeal!14,
      rounded corners=2pt,
      inner xsep=9pt,
      inner ysep=9pt
    ] {
      \begin{minipage}{0.98\textwidth}
        \centering
        {\fontsize{25}{27}\selectfont\sffamily\bfseries
          \textcolor{mDarkTeal}{Huffman $\longrightarrow{\sim}$}%
          \textcolor{mLightBrown}{4.2}%
          \hspace{0.18em}%
          \textcolor{mDarkTeal}{ bits/ch}}

        \vspace{0.75em}
        \setlength{\tabcolsep}{0pt}
        \renewcommand{\arraystretch}{1.95}
        \begin{tabular*}{\textwidth}{@{\extracolsep{\fill}}llllll}
          \codecell{a}{1111} & \codecell{b}{010101} & \codecell{c}{01011} & \codecell{d}{11100} & \codecell{e}{011} & \codecell{f}{00100} \\
          \codecell{g}{110111} & \codecell{h}{0011} & \codecell{i}{1100} & \codecell{j}{110101010} & \codecell{k}{0010110} & \codecell{l}{11101} \\
          \codecell{m}{01000} & \codecell{n}{1011} & \codecell{o}{1010} & \codecell{p}{110110} & \codecell{q}{0010111010} & \codecell{r}{1001} \\
          \codecell{s}{1000} & \codecell{t}{000} & \codecell{u}{01001} & \codecell{v}{001010} & \codecell{w}{110100} & \codecell{x}{001011100} \\
          \codecell{y}{010100} & \codecell{z}{1101010111} & \codecell{0}{00101111} & \codecell{1}{11010100} & \codecell{2}{110101110} & \codecell{3}{110101111} \\
          \codecell{4}{1101011010} & \codecell{5}{1101011110} & \codecell{6}{1101010110} & \codecell{7}{0010111011} & \codecell{8}{1101011011} & \codecell{9}{110101100}
        \end{tabular*}
      \end{minipage}
    };
  \end{tikzpicture}%
  }

  \vspace{0.2em}
  \small Frequent symbols get shorter codewords; rare symbols pay more \parencite{Huffman1952AMF}.
\end{frame}

\begin{frame}{Arithmetic Coding: Compressing ``hello''}
  \centering
  \includegraphics[width=0.94\textwidth,height=0.76\textheight,keepaspectratio]{figures/tikz/ac_interval.pdf}

  \vspace{0.2em}
  \small The message is represented by a final interval; any binary fraction inside that interval decodes to the same text \parencite{rissanenArithmeticCoding1979}.
\end{frame}

\begin{frame}{Language Modeling Objective}
  \centering
  \Large
  \[
    P(x_{1:n}) = \prod_{t=1}^{n} P(x_t \mid x_{<t})
  \]

  \vspace{0.8em}
  \[
    \mathcal{L}_{\mathrm{NTP}}
      = - \sum_{t=1}^{n} \log P_\theta(x_t \mid x_{<t})
  \]

  \vspace{0.8em}
  \normalsize
  Next-token prediction trains the same conditional probabilities an entropy coder needs.
\end{frame}

\begin{frame}{LLM as Next-Token Prediction}
  \centering
  \resizebox{\textwidth}{!}{%
  \begin{tikzpicture}[
    node distance=0.9cm and 0.9cm,
    box/.style={draw, rounded corners=2pt, minimum width=2.15cm,
                minimum height=0.85cm, align=center, font=\small},
    arr/.style={-{Stealth[length=5pt]}, thick},
  ]
    \node[box, fill=blue!12, minimum width=2.55cm] (input) {Input tokens\\\texttt{the celtics are}};
    \node[box, fill=mLightBrown!20, minimum width=2.35cm, right=of input] (model) {Language model\\e.g., GPT};
    \node[box, fill=green!14, minimum width=3.0cm, right=of model] (dist) {Predicted distribution\\
      \texttt{playing}: 0.34\\
      \texttt{favored}: 0.18\\
      \texttt{in}: 0.08\\
      \ldots};
    \node[box, fill=orange!16, minimum width=2.7cm, right=of dist] (coder) {Arithmetic encoder\\actual token = \texttt{playing}};
    \node[box, fill=red!10, minimum width=2.45cm, right=of coder] (bits) {Compressed bitstream\\\texttt{101101...}};

    \draw[arr] (input) -- (model);
    \draw[arr] (model) -- (dist);
    \draw[arr] (dist) -- (coder);
    \draw[arr] (coder) -- (bits);

    \node[font=\scriptsize, align=center] at ($(model.south)+(0,-0.55)$)
      {predicts $P(\text{token}\mid\text{context})$};
    \node[font=\scriptsize, align=center] at ($(coder.south)+(0,-0.65)$)
      {encodes the observed next token\\using the predicted distribution};
  \end{tikzpicture}
  }

\vspace{0.6em}
\small Better predictions reduce bits; every token still requires deterministic model replay.
\end{frame}

\begin{frame}{Compression Is Not Ordinary Text Generation}
  \centering
  \begin{tikzpicture}[
    box/.style={draw, rounded corners=2pt, minimum width=3.3cm,
                minimum height=0.78cm, align=center, font=\small},
    arr/.style={-{Stealth[length=5pt]}, thick},
  ]
    \node[box, fill=blue!10] (serve) at (0,1.0) {serving};
    \node[box, fill=blue!10] (sample) at (4.0,1.0) {sample token};
    \node[box, fill=blue!10] (text) at (8.0,1.0) {return text};
    \draw[arr] (serve) -- (sample);
    \draw[arr] (sample) -- (text);

    \node[box, fill=orange!14] (comp) at (0,-1.0) {compression};
    \node[box, fill=orange!14] (scores) at (4.0,-1.0) {capture scores};
    \node[box, fill=orange!14] (ac) at (8.0,-1.0) {update coder};
    \draw[arr] (comp) -- (scores);
    \draw[arr] (scores) -- (ac);
  \end{tikzpicture}

  \vspace{0.8em}
  \small Generation engines optimize sampled-token serving; compression needs ordered scores or intervals for exact encode/decode replay \parencite{kwonEfficientMemoryManagement2023,zhengSGLangEfficientExecution2024}.
\end{frame}
